
hello

input distrivution isotropic gaussian 20 dim

error comparable the optimal least squares estimator

Q: why is this suggesting effective learning algorithm?

two types of shifts, what shifts could we use?

our inputs, how many dims? here they use 20 dim inputs. 

given black sholes we will not be able to change dimensionality right?

def of learnable, do we have it in notes?

if one forawrd pass, as compared to backpropagation or what? 

squared error as loss function for training
we do not fine tune a pretrained model, we train from scratch. we do not train on actual text

can we use the three other learning algorithms for black sholes as well?

figure two should be one we can reproduce, first for linear functions and then for balck sholes derived data. or am I misunderstandiing this?

testing volatility beound the seen interval or within the seen interval but not in it. 

drawn from a different distribution on test than train, does that mean changing the sigmas or changing more in the function definition, do we need to find more degrees of freedom to introduce?

we see indicators of in-context learning, as the model error degrades close to that of the least squared estimator. This is interpreted as in-context ability extrapolating beyond the training distribution. 

the idea of trainng with and without curriculum, i dont really get it. what is the curriculum and how is it used?

what is the meaning of an unrestricted model, what is unrestricted in this context?

learning-to-learn abilities

